% ===========================================================================
% Supplementary Appendix
% GRPO for Agentic LLM Fine-Tuning
% ===========================================================================

\documentclass{article}
\usepackage[utf8]{inputenc}
\usepackage[T1]{fontenc}
\usepackage{times}
\usepackage{url}
\usepackage{latexsym}
\usepackage{graphicx}
\usepackage{amsmath,amssymb}
\usepackage{booktabs}
\usepackage{microtype}
\usepackage{hyperref}
\usepackage{caption}
\usepackage{natbib}

% ===========================================================================
% Supplementary Material
% ===========================================================================
\begin{document}

% ===========================================================================
\section{Complete Training Run Registry}
% ===========================================================================

\subsection{Tool-Use Experiments (Tinker Cloud GPU, Qwen3-8B)}

\begin{table}[htbp]
    \centering
    \caption{Complete tool-use GRPO runs on Tinker (Qwen3-8B, LoRA rank 32). All runs use Adam optimizer ($\beta_1$=0.9, $\beta_2$=0.95, $\epsilon$=1e-8), max sequence 1024 input / 512 generation.}
    \begin{tabular}{llcccccl}
        \toprule
        Run ID & Config & Group & Temp & Steps & Last-10 Reward & Status \\
        \midrule
        88ed2271 & Baseline: LR=3e-5 & 8 & 0.8 & 30 & 0.875 & Done \\
        2d488a85 & High LR: LR=1e-4 & 16 & 0.8 & 30 & 0.999 & Done \\
        5569c1fa & Low Temp: LR=3e-5 & 4 & 0.4 & 30 & 0.977 & Done \\
        22e9e5fd & xlam-60k: LR=3e-5 & 8 & 0.8 & 30 & 0.363 & Done \\
        386273f4 & Synthetic 5-tool & 4 & 0.8 & 100 & 0.825 & Done \\
        f63b618c & xlam-60k extended & 4 & 0.8 & 100 & 0.113 & Done \\
        1b01abef & MATH problems & 4 & 0.8 & 100 & 0.264 & Done \\
        \bottomrule
    \end{tabular}
    \label{tab:tool_runs}
\end{table}

\subsection{GSM8K Mathematical Reasoning Experiments (Tinker, Qwen3-8B)}

\begin{table}[htbp]
    \centering
    \caption{Complete GSM8K GRPO runs on Tinker. Multi-seed replication (Qwen3-8B), LoRA rank ablation, group size ablation, 4B multi-seed replication, and extended run. All metrics are training-set reward.}
    \begin{tabular}{lllccccccc}
        \toprule
        Run ID & Model & Seed & Rank & $G$ & LR & Steps & Peak & Last-10 \\
        \midrule
        \multicolumn{9}{l}{\emph{Multi-seed replication (Qwen3-8B, rank 32, G=4)}} \\
        899d909e & Qwen3-8B & 42 & 32 & 4 & 3e-5 & 50 & 62.5\% & 27.5\% \\
        5db4e965 & Qwen3-8B & 137 & 32 & 4 & 3e-5 & 50 & 62.5\% & 27.5\% \\
        aabb48cb & Qwen3-8B & 256 & 32 & 4 & 3e-5 & 50 & 62.5\% & 32.5\% \\
        99971b26 & Qwen3-8B & 512 & 32 & 4 & 3e-5 & 50 & 87.5\% & 30.0\% \\
        b3ba8df6 & Qwen3-8B & 999 & 32 & 4 & 3e-5 & 50 & 87.5\% & 35.0\% \\
        \midrule
        \multicolumn{9}{l}{\emph{LoRA rank ablation (Qwen3-8B, seed 42, G=4)}} \\
        ba2a1694 & Qwen3-8B & 42 & 8 & 4 & 3e-5 & 50 & 62.5\% & 21.2\% \\
        92ebcc48 & Qwen3-8B & 42 & 16 & 4 & 3e-5 & 50 & 75.0\% & 18.8\% \\
        9219771c & Qwen3-8B & 42 & 64 & 4 & 3e-5 & 50 & 87.5\% & 25.0\% \\
        \midrule
        \multicolumn{9}{l}{\emph{Group size ablation (Qwen3-8B, seed 42, rank 32)}} \\
        a13b85f7 & Qwen3-8B & 42 & 32 & 4 & 3e-5 & 50 & 62.5\% & 23.8\% \\
        8c42a6aa & Qwen3-8B & 42 & 32 & 8 & 3e-5 & 50 & 68.8\% & 24.4\% \\
        ae5c6287 & Qwen3-8B & 42 & 32 & 16 & 3e-5 & 50 & 75.0\% & 36.2\% \\
        6e2f0d73 & Qwen3-8B & 42 & 32 & 32 & 3e-5 & 50 & 100\% & 54.7\% \\
        \midrule
        \multicolumn{9}{l}{\emph{4B multi-seed replication (Qwen3.5-4B, rank 32, G=4)}} \\
        b6e8e70d & Qwen3.5-4B & 42 & 32 & 4 & 3e-5 & 50 & 87.5\% & 68.8\% \\
        566747c0 & Qwen3.5-4B & 137 & 32 & 4 & 3e-5 & 50 & 100\% & 82.5\% \\
        a00e78da & Qwen3.5-4B & 256 & 32 & 4 & 3e-5 & 50 & 100\% & 96.2\% \\
        cd17f565 & Qwen3.5-4B & 512 & 32 & 4 & 3e-5 & 50 & 100\% & 91.2\% \\
        \midrule
        \multicolumn{9}{l}{\emph{3B group-size control (Llama-3.2-3B, rank 32)}} \\
        86162fb1 & Llama-3.2-3B & 42 & 32 & 32 & 3e-5 & 50 & 12.5\% & 5.0\% \\
        \midrule
        \multicolumn{9}{l}{\emph{Extended run}} \\
        1cc20cec & Qwen3-8B & 42 & 32 & 4 & 5e-6 & 100 & 75.0\% & 27.5\% \\
        \bottomrule
    \end{tabular}
    \label{tab:gsm8k_runs}
\end{table}

\subsection{10x Structural Ceiling Experiments (Tinker, 32 runs, \$65)}

\begin{table}[htbp]
    \centering
    \footnotesize
    \caption{Complete 10x Structural Ceiling run registry. All runs are 50 steps on Tinker with full checkpoint resumption. ZVF = Zero-Variance Fraction, GU = Gradient Utilization = 1 - ZVF.}
    \begin{tabular}{llcccccc}
        \toprule
        Dimension & Config & Model & Steps & Final Reward & Avg Last-10 & Mean ZVF \\
        \midrule
        \multicolumn{7}{l}{\emph{Benchmark hierarchy (Qwen3-8B, G=16)}} \\
        Benchmark & Tool-use JSON & Qwen3-8B & 50 & 1.000 & 1.000 & -- \\
        Benchmark & GSM8K (LR=1e-4) & Qwen3-8B & 50 & 1.000 & 1.000 & $\sim$1.00 \\
        Benchmark & GSM8K (seed4) & Qwen3-8B & 50 & 0.984 & 0.972 & 0.550 \\
        Benchmark & GSM8K (seed5) & Qwen3-8B & 50 & 0.922 & 0.925 & 0.430 \\
        Benchmark & MATH-500 & Qwen3-8B & 50 & 0.720 & 0.574 & -- \\
        Benchmark & HumanEval & Qwen3-8B & 50 & 0.000 & 0.024 & -- \\
        \midrule
        \multicolumn{7}{l}{\emph{Cross-family architecture (8B, tool-use + GSM8K)}} \\
        Architecture & Tool-use & Qwen3-8B & 50 & 1.000 & 1.000 & -- \\
        Architecture & Tool-use & Llama-3.1-8B-Inst & 50 & 0.103 & -- & 0.00 \\
        Architecture & Tool-use & Llama-3.1-8B base & 50 & 0.000 & -- & 1.00 \\
        Architecture & GSM8K & Llama-3.1-8B-Inst & 39 & 0.969 & -- & 0.50 \\
        Architecture & GSM8K & Llama-3.1-8B base & 26 & 0.047 & -- & 0.75 \\
        \midrule
        \multicolumn{7}{l}{\emph{Model size ladder (GSM8K)}} \\
        Size & 8B & Qwen3-8B & 50 & 1.000 & 0.972 & 0.550 \\
        Size & 1.7B & Qwen3-1.7B & 50 & 0.016 & 0.009 & 0.885 \\
        Size & 0.6B & Qwen3-0.6B & 50 & 0.016 & 0.009 & 0.920 \\
        Size & 3B & Llama-3.2-3B & 47 & 0.016 & $\sim$0.02 & $\sim$0.90 \\
        Size & 1B & Llama-3.2-1B & 50 & 0.000 & 0.000 & 1.00 \\
        \midrule
        \multicolumn{7}{l}{\emph{Group size ablation (Qwen3-8B, GSM8K)}} \\
        Group & G=4 & Qwen3-8B & 50 & 1.000 & 0.944 & 0.520 \\
        Group & G=16 (seed4) & Qwen3-8B & 50 & 0.984 & 0.972 & 0.550 \\
        Group & G=16 (seed5) & Qwen3-8B & 50 & 0.922 & 0.925 & 0.430 \\
        Group & G=32 & Qwen3-8B & 50 & 1.000 & 0.957 & 0.455 \\
        Group & G=64 & Qwen3-8B & 50 & 1.000 & $\sim$0.98 & 0.525 \\
        \midrule
        \multicolumn{7}{l}{\emph{Learning rate ablation (Qwen3-8B, G=16, GSM8K)}} \\
        LR & 1e-5 & Qwen3-8B & 50 & 0.594 & 0.677 & 0.175 \\
        LR & 4e-5 (seed4) & Qwen3-8B & 50 & 0.984 & 0.972 & 0.550 \\
        LR & 4e-5 (seed5) & Qwen3-8B & 50 & 0.922 & 0.925 & 0.430 \\
        LR & 1e-4 & Qwen3-8B & 50 & 1.000 & 1.000 & $\sim$1.00 \\
        LR & 3e-4 & Qwen3-8B & 50 & 0.984 & 0.901 & 0.565 \\
        \midrule
        \multicolumn{7}{l}{\emph{Constrained decoding (Qwen3-8B, tool-use)}} \\
        Constrained & Unconstrained & Qwen3-8B & 50 & 0.998 & -- & 0.725 \\
        Constrained & Constrained & Qwen3-8B & 50 & 0.981 & -- & 0.660 \\
        \bottomrule
    \end{tabular}
    \label{tab:10x_runs}
\end{table}

\subsection{Google Colab QLoRA Experiments}

\begin{table}[htbp]
    \centering
    \caption{Local QLoRA experiments on Google Colab (T4 GPU, 16GB VRAM).}
    \begin{tabular}{lllccl}
        \toprule
        Model & Dataset & Method & LoRA Rank & Steps & Result \\
        \midrule
        Qwen2.5-0.5B & 14 synthetic & SFT only & 16 & 10 & JSON 30\% \\
        Qwen2.5-1.5B & Glaive-v2 & SFT+GRPO & 16 & 700 & JSON 92\% \\
        Qwen2.5-3B & ToolBench-v1 & SFT+GRPO & 8 & 40 & 0.91 multi-turn \\
        Llama-3.2-3B & GSM8K & GRPO only & 32 & 50 & 2.34\% acc \\
        Qwen3-4B & GSM8K & SFT+GRPO & 32 & 50 & In progress \\
        \bottomrule
    \end{tabular}
    \label{tab:colab_runs}
\end{table}

\subsection{Advanced RL Algorithm Variants (Colab Notebook)}

We provide a reproducible Colab notebook (\texttt{advanced\_rl\_colab.ipynb}) implementing three non-standard RL algorithms for comparison with standard GRPO:

\begin{table}[htbp]
    \centering
    \caption{Advanced RL algorithm configurations (Qwen3-8B + LoRA rank 32 on GSM8K).}
    \begin{tabular}{llcccc}
        \toprule
        Algorithm & Key Change & $\beta$ (KL) & LR & Temp & Steps \\
        \midrule
        Standard GRPO & Baseline & 0.04 & 4e-5 & 1.0 & 50 \\
        Dr.\ GRPO & No KL + entropy bonus & 0.0 & 5e-5 & 1.0 & 50 \\
        DAPO & Overlong filter + high temp & 0.0 & 4e-5 & 1.2 & 50 \\
        DPO & Offline preference pairs & 0.1 & 5e-6 & -- & 3 epochs \\
        \bottomrule
    \end{tabular}
    \label{tab:advanced_rl}
\end{table}

\noindent\textbf{Dr.\ GRPO} \citep{deepseekr1} removes the KL divergence penalty and adds a token-level entropy bonus ($\lambda=0.01$) to prevent the entropy collapse documented in Section~5.3 of the main paper. \textbf{DAPO} \citep{yu2025dapo} further decouples alignment by using asymmetric advantage clipping and higher sampling temperature for diverse rollouts. \textbf{DPO} \citep{rafailov2024dpo} uses rejection-sampled preference pairs (correct vs.\ incorrect completions from the base model) for offline preference learning without on-policy rollouts.

% ===========================================================================
\section{Claim-to-Configuration Mapping}
% ===========================================================================

Table~\ref{tab:config_map} maps each main-paper claim to its exact training configuration, enabling reviewers to trace results without cross-referencing run IDs.

\begin{table}[htbp]
    \centering
    \footnotesize
    \caption{Configuration for each main-paper claim. $N$ = number of independent runs or seeds.}
    \begin{tabular}{p{3.2cm}lcccccc}
        \toprule
        Claim & Model & Rank & $G$ & LR & Steps & $N$ \\
        \midrule
        \multicolumn{7}{l}{\emph{Original experiments}} \\
        Tool 0\%$\to$92\% & Qwen2.5-1.5B & 16 & 8 & 3e-5 & 700 & 1 \\
        Multi-turn 0.72$\to$0.91 & Qwen2.5-3B & 8 & 4 & 3e-5 & 40 & 1 \\
        Code 32\%$\to$40\% & Qwen3-8B & 32 & 4 & 3e-5 & 50 & 1 \\
        GSM8K 30.5\% (SD 3.3) & Qwen3-8B & 32 & 4 & 3e-5 & 50 & 5 \\
        Held-out 83.3\% & Qwen3-8B & 32 & 4 & 3e-5 & 50 & 5 \\
        Group-size scaling & Qwen3-8B & 32 & 4--32 & 3e-5 & 50 & 4 \\
        Capacity threshold & 3B/4B & 32 & 4,32 & 3e-5 & 50 & 2/4 \\
        MoE volatility & Qwen3-30B & 32 & 8 & 3e-5 & 50 & 1 \\
        LoRA rank tradeoff & Qwen3-8B & 8--64 & 4 & 3e-5 & 50 & 3 \\
        \midrule
        \multicolumn{7}{l}{\emph{10x Structural Ceiling (32 runs)}} \\
        Benchmark hierarchy & Qwen3-8B & 32 & 16 & 4e-5 & 50 & 6 \\
        Architecture depend. & Qwen/Llama-8B & 32 & 16 & 4e-5 & 50 & 5 \\
        Size ladder & 0.6B--8B & 32 & 16 & 4e-5 & 50 & 5 \\
        Group saturation & Qwen3-8B & 32 & 4--64 & 4e-5 & 50 & 5 \\
        LR tradeoff & Qwen3-8B & 32 & 16 & 1e-5--3e-4 & 50 & 5 \\
        Constrained decoding & Qwen3-8B & 32 & 16 & 4e-5 & 50 & 2 \\
        SFT prerequisite & Llama-8B base/inst & 32 & 16 & 4e-5 & 50 & 2 \\
        Reward hacking & Llama-8B base & 32 & 16 & 4e-5 & 50 & 1 \\
        \bottomrule
    \end{tabular}
    \label{tab:config_map}
\end{table}

% ===========================================================================
\section{Claim-to-Run Mapping}
% ===========================================================================

Table~\ref{tab:claim_map} maps each main-paper finding to the specific runs that support it, enabling reviewers to trace every claim to its evidence.

\begin{table}[htbp]
    \centering
    \footnotesize
    \caption{Mapping of main-paper claims to supporting evidence. Single-run claims are marked ($\star$).}
    \begin{tabular}{p{3.8cm}lp{4.5cm}}
        \toprule
        Claim & Runs & Evidence \\
        \midrule
        \multicolumn{3}{l}{\emph{Original experiments}} \\
        Tool calling 0\%$\rightarrow$92\% & Colab Qwen2.5-1.5B & Single SFT+GRPO pipeline ($\star$) \\
        Multi-turn 0.72$\rightarrow$0.91 & Colab Qwen2.5-3B & 4 scenarios, custom rubric ($\star$) \\
        Code 32\%$\rightarrow$40\% & Tinker Qwen3-8B & 50-item subset, $p{=}0.53$ ($\star$) \\
        GSM8K 30.5\% (SD 3.3\%) & 899d, 5db4, aabb, 9997, b3ba & 5 seeds, same config \\
        Held-out 83.3\% (SD 2.2\%) & Same 5 seeds & 200 test examples each \\
        Capacity threshold 3B--4B & Colab 3B + 8616 + 4$\times$4B & 3B confirmed at both $G$; 4 4B seeds \\
        MoE volatility 2.43$\times$ & MoE run + 5 dense seeds & Levene's $p{=}7{\times}10^{-6}$ \\
        SFT+GRPO complementarity & Colab 1.5B, 3B & 2 pipelines ($\star$) \\
        Group-size scaling & a13b, 8c42, ae5c, 6e2f & 4 group sizes, seed 42 \\
        LoRA rank--speed tradeoff & ba2a, 92eb, 9219 & 3 ranks, seed 42 ($\star$) \\
        Two-phase learning & Arithmetic diagnostic run & Single run ($\star$) \\
        Synthetic--real gap & 386273f4, f63b618c & 2 runs, same model \\
        \midrule
        \multicolumn{3}{l}{\emph{10x Structural Ceiling (32 runs, W\&B: tinker-structural-ceiling)}} \\
        Benchmark hierarchy & 6 runs across 4 benchmarks & Tool/GSM8K/MATH/HumanEval \\
        Architecture dependence & 5 runs (Qwen/Llama-8B) & Tool-use + GSM8K, both families \\
        SFT prerequisite & 2 runs (base vs instruct) & Llama-8B, +0.922 delta \\
        Size ladder & 5 runs (0.6B--8B) & Qwen + Llama, cross-family \\
        Group saturation (ZVF/GU) & 5 runs (G=4/16/32/64) & Novel diagnostic metric \\
        LR speed--saturation & 5 runs (1e-5 to 3e-4) & Full 50-step correction \\
        Constrained decoding & 2 runs & No difference (0.998 vs 0.981) \\
        Reward hacking/collapse & 1 run (Llama-8B base) & Loss $-238$, collapse at step 41 ($\star$) \\
        \bottomrule
    \end{tabular}
    \label{tab:claim_map}
\end{table}

\textbf{Key limitation:} Of the original twelve claims, four rest on single runs or single pipelines (marked $\star$). The 10x Structural Ceiling experiment adds 32 systematic runs across 7 dimensions, substantially strengthening the evidence base. The capacity threshold is now confirmed across both Qwen (0.6B, 1.7B) and Llama (1B, 3B) families, partially resolving the architecture confound. The reward hacking observation remains a single run.

% ===========================================================================
\section{Training Hyperparameters}
% ===========================================================================

\subsection{Tinker Cloud GPU Configuration}

\begin{table}[htbp]
    \centering
    \caption{Tinker SDK hyperparameters for all GRPO runs.}
    \begin{tabular}{ll}
        \toprule
        Parameter & Value \\
        \midrule
        SDK Version & v0.16.1 \\
        Optimizer & Adam \\
        $\beta_1$ & 0.9 \\
        $\beta_2$ & 0.95 \\
        $\epsilon$ & 1e-8 \\
        Weight Decay & 0 \\
        Batch Size & 2 prompts/step \\
        Group Size & 4, 8, 16, or 32 \\
        Max Input Length & 1024 tokens \\
        Max Generation Length & 512 tokens \\
        Warmup & None (constant LR) \\
        Checkpoint Frequency & Final only \\
        \bottomrule
    \end{tabular}
    \label{tab:tinker_hyperparams}
\end{table}

\subsection{LoRA Configuration}

\begin{table}[htbp]
    \centering
    \caption{LoRA adapter targets and ranks used across experiments.}
    \begin{tabular}{ll}
        \toprule
        Parameter & Value \\
        \midrule
        Target Modules & q\_proj, k\_proj, v\_proj, o\_proj, \\
                        & gate\_proj, up\_proj, down\_proj \\
        LoRA Ranks Tested & 8, 16, 32, 64 \\
        Default Rank & 32 \\
        Alpha & 2 $\times$ rank (default) \\
        Dropout & 0.0 \\
        QLoRA (Colab) & 4-bit NF4, bfloat16 compute \\
        \bottomrule
    \end{tabular}
    \label{tab:lora_config}
\end{table}

% ===========================================================================
\section{Failure Case Analysis}
% ===========================================================================

\subsection{Three-Billion Parameter Failure Mode}

The Llama-3.2-3B model on GSM8K achieved only 2.34\% accuracy after 50 GRPO steps (95\% CI $[0.3\%, 4.4\%]$ via binomial proportion on $G{=}4$ groups), barely above the 1.56\% random baseline. A one-sample binomial test yields $p{=}0.68$, failing to reject the null hypothesis that the model performs at chance. This constitutes a \textbf{formal negative result}: GRPO at $G{=}4$ does not produce statistically detectable learning in 3B models on GSM8K.

The root cause is that \textbf{56\% of training steps had zero gradient} because all $G{=}4$ sampled completions received identical rewards (all incorrect). When all completions are equally wrong, GRPO advantages are zero and no gradient update occurs. The 8B model (Qwen3-8B) shows only 16--28\% zero-loss rate, enabling effective gradient updates.

\textbf{Practitioner guideline:} Sub-3B dense models should not be expected to learn via GRPO with $G \leq 4$ on math reasoning tasks. Higher group sizes ($G{=}16$--$64$) or curriculum warm-starting may be needed to generate within-group reward variance at this scale, but this remains untested.

\begin{table}[htbp]
    \centering
    \caption{Zero-loss step percentage across model sizes (lower is better).}
    \begin{tabular}{lccc}
        \toprule
        Model & Parameters & Zero-Loss \% & Final Acc \\
        \midrule
        Llama-3.2-3B & 3B & 56\% & 2.34\% \\
        Qwen3-8B (seed 512) & 8B & 16\% & 30.0\% \\
        Qwen3-8B (seed 137) & 8B & 28\% & 27.5\% \\
        Qwen3-8B (seed 256) & 8B & 24\% & 32.5\% \\
        \bottomrule
    \end{tabular}
    \label{tab:zero_loss}
\end{table}

\subsection{Comprehensive Zero-Loss Dynamics Across All GSM8K Runs}

Table~\ref{tab:zero_loss_comprehensive} reports zero-loss and zero-reward step frequencies across all available GSM8K training runs, providing the quantitative basis for the capacity-threshold and exploration-sparsity claims in the main paper.

\begin{table}[htbp]
    \centering
    \caption{Zero-loss and zero-reward dynamics across all GSM8K GRPO runs. Zero-loss indicates all-same-reward groups (no gradient signal). Zero-reward indicates all-incorrect completions. 3B data from Colab; 4B and 8B from Tinker.}
    \begin{tabular}{llccccc}
        \toprule
        Model & Config & Steps & Zero-Loss \% & Zero-Reward \% & Peak Acc & Last-10 Acc \\
        \midrule
        Llama-3.2-3B & seed=42, G=4 & 50 & 56\% & 56\% & 3.1\% & 2.3\% \\
        Llama-3.2-3B & seed=42, G=32 & 50 & 18\% & 18\% & 12.5\% & 5.0\% \\
        \midrule
        Qwen3.5-4B & seed=42, G=4 & 50 & -- & -- & 87.5\% & 68.8\% \\
        Qwen3.5-4B & seed=137, G=4 & 50 & 68\% & 0\% & 100\% & 82.5\% \\
        Qwen3.5-4B & seed=256, G=4 & 50 & -- & -- & 100\% & 96.2\% \\
        Qwen3.5-4B & seed=512, G=4 & 50 & -- & -- & 100\% & 91.2\% \\
        \midrule
        Qwen3-8B & seed=137, r=32, G=4 & 50 & 28\% & 24\% & 62.5\% & 27.5\% \\
        Qwen3-8B & seed=256, r=32, G=4 & 50 & 24\% & 24\% & 62.5\% & 32.5\% \\
        Qwen3-8B & seed=512, r=32, G=4 & 50 & 16\% & 10\% & 87.5\% & 30.0\% \\
        Qwen3-8B & seed=42, r=32, G=4 & 50 & 18\% & 14\% & 62.5\% & 27.5\% \\
        Qwen3-8B & seed=999, r=32, G=4 & 50 & 18\% & 16\% & 87.5\% & 35.0\% \\
        Qwen3-8B & seed=42, r=8, G=4 & 50 & 20\% & 14\% & 62.5\% & 21.2\% \\
        Qwen3-8B & seed=42, r=16, G=4 & 50 & 20\% & 18\% & 75.0\% & 18.8\% \\
        Qwen3-8B & seed=42, r=64, G=4 & 50 & 18\% & 16\% & 87.5\% & 25.0\% \\
        \midrule
        \multicolumn{7}{l}{\emph{Group size ablation (Qwen3-8B, seed 42, rank 32)}} \\
        Qwen3-8B & seed=42, G=4 & 50 & 26\% & 20\% & 62.5\% & 23.8\% \\
        Qwen3-8B & seed=42, G=8 & 50 & 6\% & 6\% & 68.8\% & 24.4\% \\
        Qwen3-8B & seed=42, G=16 & 50 & 2\% & 2\% & 75.0\% & 36.2\% \\
        Qwen3-8B & seed=42, G=32 & 50 & 2\% & 0\% & 100\% & 54.7\% \\
        \midrule
        Qwen3-8B & lr=5e-6, G=4, 100 steps & 100 & 17\% & 15\% & 75.0\% & 27.5\% \\
        \bottomrule
    \end{tabular}
    \label{tab:zero_loss_comprehensive}
\end{table}

\textbf{Key patterns:}
\begin{itemize}
    \item The 3B model's 56\% zero-loss rate stems from zero-reward (all-incorrect) groups, confirming reward sparsity as the proximate cause.
    \item The 4B model's 68\% zero-loss rate stems from \emph{all-correct} groups (0\% zero-reward), indicating the opposite regime: the model solves problems so easily that group variance vanishes. This is productive zero-loss (learning saturated, not stalled).
    \item 8B models show 16--28\% zero-loss, with seed 512 showing the lowest rate and highest peak accuracy, suggesting exploration varies meaningfully across seeds.
    \item LoRA rank does not substantially affect zero-loss rates (18--20\% for r=8/16/64), supporting the paper's finding that rank affects learning speed but not the fundamental capacity to generate reward signal.
    \item \textbf{Group size dramatically reduces zero-loss:} At $G{=}4$, 26\% of steps have zero gradient; at $G{=}32$, only 2\%. Larger groups ensure at least one correct and one incorrect completion per group, providing consistent gradient signal. This explains why $G{=}32$ achieves 54.7\% vs.\ $G{=}4$'s 23.8\%.
\end{itemize}

\subsection{Qwen3-30B-MoE Volatility Analysis}

The Qwen3-30B-MoE model reached 99\% peak GSM8K training-step accuracy in our internal runs but exhibited 2.43$\times$ higher step-to-step reward variance than the dense 8B model. Statistical testing (Levene's test) confirms this difference is significant: $F(1, 98) = 22.4$, $p = 7.0 \times 10^{-6}$.

This volatility stems from MoE routing: different experts activate for different tokens, creating non-smooth loss surfaces. The sparse activation pattern means small changes in the router can dramatically shift which experts process each token, amplifying training instability.

\begin{table}[htbp]
    \centering
    \caption{MoE vs. Dense training stability comparison.}
    \begin{tabular}{lccc}
        \toprule
        Model & Step Var (SD) & Peak Acc & Final Acc \\
        \midrule
        Qwen3-8B Dense & 0.12 & 87.5\% & 30.0\% \\
        Qwen3-30B-MoE & 0.29 & 99.0\% & 95.0\% \\
        Ratio (MoE/Dense) & 2.43$\times$ & --- & --- \\
        \bottomrule
    \end{tabular}
    \label{tab:moe_volatility}
\end{table}

\subsection{Policy Entropy, KL, and Group-Composition Diagnostics}
\label{sec:diagnostics_table}

The arithmetic task provides the richest telemetry logged during GRPO training. Table~\ref{tab:diagnostics} summarizes key diagnostic metrics at milestone training steps. Policy entropy collapses by $>$3 orders of magnitude (0.773 $\rightarrow$ 0.0002) within 12 steps, while KL to the training policy remains near zero throughout, indicating the model converges to a near-deterministic policy without drifting far from initialization. Group composition shifts from 81\% mixed groups at step 0 (productive gradient signal) to $>$96\% all-good groups by step 4 (saturated learning, diminishing returns).

\begin{table}[htbp]
    \centering
    \caption{Diagnostic telemetry from arithmetic GRPO training (Qwen3-8B, 100 steps, G=4). Entropy, KL, and group composition are logged per step; selected milestones shown.}
    \begin{tabular}{rcccccc}
        \toprule
        Step & Reward & Entropy & KL (v1) & Frac Mixed & Frac All Good & Frac All Bad \\
        \midrule
        0  & 0.676 & 0.773   & 5.7e-4  & 0.81 & 0.19 & 0.00 \\
        1  & 0.795 & 0.586   & 6.3e-4  & 0.62 & 0.37 & 0.01 \\
        2  & 0.933 & 0.228   & $-$1.4e-4 & 0.36 & 0.64 & 0.00 \\
        4  & 0.999 & 0.030   & 9.7e-5  & 0.04 & 0.96 & 0.00 \\
        7  & 0.990 & 0.004   & 3.0e-5  & 0.00 & 0.99 & 0.01 \\
        11 & 1.000 & 2.5e-4  & $-$6.6e-6 & 0.00 & 1.00 & 0.00 \\
        19 & 0.998 & 0.006   & $-$4.5e-5 & 0.01 & 0.99 & 0.00 \\
        \bottomrule
    \end{tabular}
    \label{tab:diagnostics}
\end{table}

\textbf{Key observations:}
\begin{itemize}
    \item \textbf{Entropy collapse:} 0.773 $\rightarrow$ 0.0002 (3,865$\times$ reduction) in $\sim$12 steps. Without explicit entropy regularization, the policy becomes near-deterministic, which limits further exploration.
    \item \textbf{KL remains near zero:} KL divergence to the initial training policy oscillates around $\pm$0.001 nats, indicating no KL anchor was needed for this short horizon / easy task. Longer or harder tasks may require explicit KL regularization.
    \item \textbf{Group saturation:} Once frac\_all\_good exceeds 0.95, within-group variance drops, advantages shrink, and effective learning rate diminishes (the ``zero-advantage'' regime). This provides a natural stopping criterion.
    \item \textbf{Missing for GSM8K/MATH:} Comparable entropy/KL/group-composition metrics were \emph{not logged} for the main GSM8K and MATH experiments (Tinker SDK did not expose these in the training loop). Adding this telemetry is a Phase 1 priority in our experimental roadmap (Section 7 of the main paper).
\end{itemize}

\subsection{MBPP Evaluation Failure}

The MBPP (Mostly Basic Python Problems) benchmark produced 0\% pass rate across all runs due to an output parsing bug in the evaluation script. The script expected a specific JSON format for extracted answers but the model output format differed, preventing correct answer matching.

This represents a methodological failure: an invalidated benchmark should not be reported. Only HumanEval results (which used a working evaluator) are reported in the main paper. Future work should verify evaluator correctness before running experiments.

% ===========================================================================
\section{Additional Experimental Details}
% ===========================================================================

\subsection{Reward Function Weighting Rationale}

The tool calling reward (0.3 + 0.4 + 0.3) was designed to prioritize:
\begin{itemize}
    \item \textbf{Correct tool name (+0.4):} Selecting the right tool is the core judgment decision
    \item \textbf{Valid JSON (+0.3):} Required for downstream parsing
    \item \textbf{Arguments present (+0.3):} Functional tool use requires populated arguments
\end{itemize}

These weights were not validated via ablation. Future work should test alternative weightings, particularly increasing the tool-name weight to further emphasize selection judgment over format compliance.

\subsection{Evaluation Protocol Limitations}

As noted in the main paper, our primary evaluation metrics are computed on \textbf{training prompts with new rollouts}, not on held-out test sets. Specifically:
\begin{itemize}
    \item Tool-use rewards: Fresh sampling from the same prompt distribution
    \item GSM8K accuracy: Random subset of 8 training examples per step (batch=2, group=4)
    \item HumanEval: Standard pass@1 on the 50-example benchmark (test set)
\end{itemize}

The GSM8K training-step evaluation is the most compromised: evaluating on training data with small group sizes (G=4) means the ``test'' examples overlap with training distribution. We subsequently performed held-out evaluation on 200 randomly sampled GSM8K test-set examples (5 seeds, 200 each), reporting 83.3\% (SD 2.2\%) accuracy. A base-model control (Qwen3-8B without LoRA) scored 82.0\% on the same 200 examples, showing the +1.3pp delta is not significant ($p{=}0.26$).

\subsection{Multi-Turn Tool Calling Reward Breakdown}

The multi-turn reward function (5-component) enables finer-grained learning signals:

\begin{table}[htbp]
    \centering
    \caption{Multi-turn tool calling reward components.}
    \begin{tabular}{ll}
        \toprule
        Component & Weight \\
        \midrule
        First turn = tool call & +0.25 \\
        Final natural language answer & +0.30 \\
        Arguments populated & +0.15 \\
        Clean JSON output & +0.10 \\
        Repeated tool call (penalty) & -0.30 \\
        \bottomrule
    \end{tabular}
    \label{tab:multiturn_reward}
\end{table}

The -0.30 penalty for repeated tool calls was critical: it eliminated SFT's looping failure mode entirely (Section 4.1.2 of main paper). Without this penalty, models would repeatedly call the same tool in an infinite loop.

Because these scores come from custom scenario rollouts rather than a standard benchmark, we emphasize that they are internal protocol numbers. We did not measure inter-rater reliability, did not use multiple human judges, and did not adopt proxy-state final-answer evaluation. Future revisions should benchmark the same agents under ToolRM / FC-RewardBench-style outcome evaluation or other standardized multi-turn frameworks.

% ===========================================================================
\section{Qualitative Examples}
% ===========================================================================

\subsection{Successful Tool Call (Post-GRPO)}
\begin{verbatim}
<|im_start|>user
What is the weather in Bangalore?<|im_end|>
<|im_start|>assistant
{
  "tool_calls": [
    {
      "name": "get_weather",
      "arguments": {
        "location": "Bangalore",
        "unit": "celsius"
      }
    }
  ]
}<|im_end|>
\end{verbatim}

\subsection{Failed Tool Call (Pre-GRPO, Plain Text)}
\begin{verbatim}
<|im_start|>user
What is the weather in Bangalore?<|im_end|>
<|im_start|>assistant
The weather in Bangalore is sunny and warm today.
The temperature is around 28 degrees Celsius.
It should be a nice day to go outside!<|im_end|>
\end{verbatim}

The pre-GRPO model generates natural language responses instead of structured tool calls. After GRPO training, the model reliably produces valid JSON with correct tool names and populated arguments.

\subsection{Two-Phase Learning Trajectory (Representative)}

\begin{verbatim}
Step 1:  Reward=0.00, Loss=0.85, Output="42" (format: plain text)
Step 5:  Reward=0.05, Loss=0.80, Output="The answer is 42." (format: text with number)
Step 10: Reward=0.10, Loss=0.72, Output="\boxed{42}" (format: boxed answer)
Step 15: Reward=0.14, Loss=0.65, Output="\boxed{42}" (stable format)
Step 20: Reward=0.40, Loss=0.55, Output="\boxed{42}" (reasoning begins)
Step 25: Reward=0.58, Loss=0.40, Output="\boxed{42}" (reasoning improves)
Step 30: Reward=0.55, Loss=0.38, Output="\boxed{42}" (plateau)
\end{verbatim}

This trajectory illustrates the two-phase pattern: format compliance (steps 1-20) precedes reasoning improvement (steps 20-30).

% ===========================================================================
\section{Platform and Compute Costs}
% ===========================================================================

\begin{table}[htbp]
    \centering
    \caption{Compute platform usage and estimated costs.}
    \begin{tabular}{llcl}
        \toprule
        Platform & GPU & Experiments & Est. Cost \\
        \midrule
        Tinker SDK v0.16.1 & Cloud A100 & 25 original runs & \$40--55 \\
        Tinker SDK v0.16.1 & Cloud A100 & 32 runs (10x Structural Ceiling) & $\sim$\$65 \\
        Google Colab Pro & T4 16GB & 5 team members & \$10/person \\
        HuggingFace Hub & N/A & Model hosting & Free \\
        \bottomrule
    \end{tabular}
    \label{tab:compute}
\end{table}

Total Tinker spend: $\sim$\$105--120 across 57 runs. Total estimated project spend: $\sim$\$130, well within typical capstone project budgets.

\begin{table}[htbp]
    \centering
    \caption{Estimated per-run compute budgets for representative experiments. Token estimates assume batch\_size=2, group\_size as configured, max\_generation=512 tokens. Wall-clock times are from logged telemetry where available; others are estimated from Tinker billing.}
    \begin{tabular}{llcccc}
        \toprule
        Domain & Config & Steps & Group & Est. Tokens & Est. Time \\
        \midrule
        Arithmetic & Qwen3-8B, r=32 & 100 & 4 & 410K & $\sim$17 min \\
        GSM8K & Qwen3-8B, r=32 & 50 & 8 & 410K & $\sim$25 min \\
        GSM8K & Qwen3-8B, r=32 & 100 & 8 & 820K & $\sim$50 min \\
        GSM8K & Qwen3.5-4B, r=32 & 50 & 8 & 410K & $\sim$20 min \\
        Tool (synthetic) & Qwen3-8B, r=32 & 30 & 8 & 245K & $\sim$15 min \\
        Tool (xlam-60k) & Qwen3-8B, r=32 & 100 & 8 & 820K & $\sim$50 min \\
        \bottomrule
    \end{tabular}
    \label{tab:per_run_budget}
\end{table}

\textbf{Limitations of budget reporting:} We do not report exact tokens processed per run (the Tinker SDK does not expose this metric). Token estimates above are upper bounds assuming maximum generation length. Wall-clock times exclude checkpoint saving overhead. GPU-hours are not tracked per-run; the aggregate across all 57 Tinker runs is estimated at 24--32 A100 GPU-hours (original 25 runs: 12--16 hours; 10x Structural Ceiling 32 runs: $\sim$12--16 hours at $\sim$40 min wall time for parallel batch).

We report GPU platform and wall-clock budget, but not a full token-processed accounting per run. Likewise, our data-split story is uneven across domains: HumanEval uses a held-out benchmark subset, while tool calling and GSM8K largely report training-distribution rollouts. This split mismatch is a major threat to cross-domain comparability.

% ===========================================================================
\section{Code and Reproducibility}
% ===========================================================================

All training scripts are available in the repository:
\begin{itemize}
    \item \texttt{grpo\_gsm8k\_base.py} -- Parameterized GSM8K training
    \item \texttt{grpo\_exp\_a\_baseline.py}, \texttt{exp\_b\_high\_lr.py}, \texttt{exp\_c\_low\_temp.py}, \texttt{exp\_d\_xlam.py} -- Tool-use experiments
    \item \texttt{grpo\_100\_xlam.py}, \texttt{grpo\_100\_synthetic.py}, \texttt{grpo\_100\_math.py} -- Extended runs
    \item \texttt{experiments/10x\_structural\_ceiling/grpo\_10x\_runner.py} -- Unified runner with YAML config, checkpoint resumption, ZVF/GU logging
    \item \texttt{experiments/10x\_structural\_ceiling/configs/*.yaml} -- 23 YAML configs for all 10x ablation dimensions
    \item \texttt{experiments/10x\_structural\_ceiling/RESULTS.md} -- Complete 50-step results with 8 findings
\end{itemize}

At present, we can release training scripts, paper sources, and the hardened held-out evaluation script, but the exact prompt packs / scenario definitions used for every internal tool-calling score are not yet packaged as a replication bundle. Releasing prompts, scoring rubrics, and scenario manifests is an explicit next step for reproducibility.

For stability diagnostics, we did not log KL-to-SFT trajectories or policy entropy, and our MoE analysis does not yet include routing entropy or expert-load imbalance measurements. Those omissions limit causal interpretation of both drift and volatility claims.

Representative training logs were collected during experimentation, but the literal local \texttt{/tmp} paths used during development are not a portable reproduction interface and should not be treated as a release artifact.

Checkpoint access is partial rather than universal at present:
\begin{itemize}
    \item Some trained artifacts were published to project-managed model hosting during development.
    \item Some tool-use and math checkpoints were written to Tinker under run-specific paths such as \texttt{tinker://\textit{run\_id}/sampler\_weights/final}, but present-day recovery depends on retention and account access.
\end{itemize}
A clean anonymous release should replace personal account identifiers with project-level artifact locations or explicit recovery instructions.

% ===========================================================================
\section{Content Moved from Main Paper}
% ===========================================================================

\subsection{GRPO Hyperparameter Comparison Table}

\begin{table}[htbp]
    \centering
    \footnotesize
    \caption{Standard GRPO hyperparameters across published systems.}
    \begin{tabular}{lcccc}
        \toprule
        System & $G$ & LR & $T$ & $\beta$ (KL) \\
        \midrule
        DeepSeek-R1-Zero & 16 & $3 \times 10^{-6}$ & 1.0 & 0.001 \\
        DeepSeekMath-RL 7B & 64 & $1 \times 10^{-6}$ & -- & 0.04 \\
        JT-Math & 16 & $4 \times 10^{-6}$ & 1.2 & 0.001 \\
        Skywork-OR1 & -- & -- & -- & 0 \\
        R1-style (Qwen/Llama) & 8 & $2.5$--$5 \times 10^{-6}$ & 0.7 & 0 \\
        \midrule
        \textbf{Ours} & \textbf{16} & $\mathbf{4 \times 10^{-5}}$ & \textbf{1.0} & \textbf{--} \\
        \bottomrule
    \end{tabular}
\end{table}

Our $\sim$10$\times$ higher learning rate is attributed to LoRA-only training (0.1\% trainable parameters vs.\ full fine-tuning in published systems).

\subsection{Reward Function Limitations}

Our reward functions are deliberately simple and verifiable, but this coarseness introduces known limitations. The tool-calling reward checks JSON validity, tool name, and argument-key presence but does \emph{not} verify argument-value correctness---a model can receive full reward by calling the right function with plausible but wrong values. This creates a reward-hacking surface. Adding argument-type and value checks, or leveraging outcome reward models such as ToolRM, would reduce this risk. The math reward is binary (correct/incorrect) with no partial credit for intermediate reasoning steps, contributing to reward sparsity for weaker models.

\subsection{Safety Considerations}

GRPO-trained tool-calling models introduce safety risks: (1) \textbf{incorrect tool execution}---valid JSON with wrong functions or dangerous arguments; (2) \textbf{persistent loops}---the -0.30 penalty eliminated loops in our experiments but may not suffice for complex workflows; (3) \textbf{hallucinated capabilities}---calls to nonexistent tools, especially when transferring from synthetic to real schemas. Future work should incorporate hallucination and loop penalties validated by standardized evaluators.

\subsection{Practitioner Decision Guide}

Based on our findings, we offer preliminary guidance for practitioners considering GRPO for post-training (Table~\ref{tab:practitioner}).

\begin{table}[htbp]
    \centering
    \footnotesize
    \caption{Practitioner decision guide based on empirical findings. All recommendations are preliminary and derived from limited experiments.}
    \begin{tabular}{lll}
        \toprule
        Condition & Recommendation & Evidence \\
        \midrule
        Model $\leq$3B & Do not use GRPO ($G{=}32$ yields only 5\%) & Confirmed at $G{=}4$ and $G{=}32$ (Sec.~5.1) \\
        Model 3B--4B & Use $G \geq 16$; consider SFT warm-start & Capacity threshold (Sec.~5.1) \\
        Model $\geq$8B & GRPO viable; $G \geq 16$ recommended; rank 32 sufficient & 5-seed replication + group ablation \\
        MoE architecture & Expect 2--3$\times$ reward variance & $p{=}7{\times}10^{-6}$ (Sec.~7) \\
        Tool calling & SFT first, then GRPO & 0\%$\rightarrow$92\% (Sec.~4.1) \\
        Synthetic data & Expect 3--8$\times$ easier than real & Sec.~5.2 \\
        LoRA rank choice & Rank 32 default; 64 for faster convergence & Rank ablation (Sec.~4.4.3) \\
        \bottomrule
    \end{tabular}
    \label{tab:practitioner}
\end{table}

\textbf{Caveat:} These guidelines are derived from experiments on Qwen/Llama families with GSM8K and custom tool-calling tasks. Generalization to other model families, harder benchmarks, or multi-turn agentic trajectories is untested.

\subsection{Planned Experiments and Evaluation Roadmap}

\paragraph{Phase 1 (Immediate):} Held-out GSM8K test evaluation (1,319 examples, 3+ seeds, bootstrapped 95\% CIs) and MATH test set transfer.

\paragraph{Phase 2 (Ablations):} Group size sweep $G \in \{4, 8, 16, 32\}$ on 8B completed; remaining: $G \in \{64, 128\}$, group-size ablation on 3B, temperature sweep $T \in \{0.7, 1.0, 1.3, 1.5\}$; KL anchoring $\beta \in \{0.01, 0.05, 0.1\}$; curriculum (easy-to-hard).

\paragraph{Phase 3 (Benchmarks):} Canonical HumanEval-164/MBPP with pass@$k$; FC-RewardBench-style tool evaluation; SFT-only, PPO, DPO, RLOO baselines under matched compute. Reproducible notebooks (Dr.\ GRPO, DAPO, DPO, PPO, REINFORCE) are available in the supplementary repository.

\paragraph{Phase 4 (Diagnostics):} MoE routing telemetry; LoRA rank $r \in \{8, 16, 32, 64\}$ with memory/throughput/accuracy; QR-Adaptor comparison.

Estimated scope: $\sim$20--30 runs, \$500--\$1000.

\subsection{Reviewer Priority Analysis}

A systematic analysis of 490 RL-for-LLM reviews across 248 papers (2017--2023) finds that \emph{insufficient evaluation} is flagged in 86.3\% of reviews, with \emph{missing baselines} (56.8\%) and \emph{missing ablations} (14.5\%) as the top complaints. Human evaluation absence is rarely criticized (2.1\%), and including capability benchmarks shows no significant correlation with reviewer scores (Mann--Whitney $U$, $p > 0.05$). This motivates our prioritization of baseline comparisons and ablation coverage over human-judge studies.

% ===========================================================================
\section{Tool-Calling Evaluation Protocol}
% ===========================================================================

This section documents the custom tool-calling evaluation setup used in the main paper, addressing the reviewer concern about non-standardized evaluation.

\subsection{Tool Schema Format}

All tool schemas follow the OpenAI-compatible function-calling format:

\begin{verbatim}
{
  "type": "function",
  "function": {
    "name": "get_weather",
    "description": "Get current weather for a location",
    "parameters": {
      "type": "object",
      "properties": {
        "location": {"type": "string", "description": "City name"},
        "unit": {"type": "string", "enum": ["celsius", "fahrenheit"]}
      },
      "required": ["location"]
    }
  }
}
\end{verbatim}

\subsection{Single-Turn Evaluation Rubric}

Each tool call is scored on three binary components:

\begin{table}[htbp]
    \centering
    \caption{Single-turn tool-calling reward rubric.}
    \begin{tabular}{lcp{6cm}}
        \toprule
        Component & Weight & Criterion \\
        \midrule
        JSON Valid & +0.3 & Output parses as valid JSON \\
        Correct Tool & +0.4 & \texttt{name} field matches expected tool \\
        Arguments Present & +0.3 & All required argument keys exist \\
        \bottomrule
    \end{tabular}
\end{table}

\textbf{Limitations:} This rubric does \emph{not} verify argument values, argument types, or semantic correctness. A model can score 1.0 by producing syntactically correct JSON with plausible but wrong argument values.

\subsection{Multi-Turn Scenario Set}

Four hand-authored scenarios with deterministic expected tool sequences:

\begin{table}[htbp]
    \centering
    \caption{Multi-turn evaluation scenarios.}
    \begin{tabular}{llc}
        \toprule
        Scenario & Expected Tool Sequence & Turns \\
        \midrule
        Weather + Packing & get\_weather $\rightarrow$ suggest\_packing & 2 \\
        Stock + News & get\_stock\_price $\rightarrow$ get\_news & 2 \\
        Search + Calculate & web\_search $\rightarrow$ calculator & 2 \\
        Single Tool & get\_weather & 1 \\
        \bottomrule
    \end{tabular}
\end{table}

\subsection{Prompt Template}

\begin{verbatim}
<|im_start|>system
You are a helpful assistant with access to the following tools:
{tool_schemas_json}

When you need to use a tool, respond with a JSON object
containing "tool_calls" array. Otherwise respond normally.
<|im_end|>
<|im_start|>user
{user_query}<|im_end|>
\end{verbatim}

\subsection{Judge Details}

Scoring is fully automated (no LLM judge, no human annotation):
\begin{itemize}
    \item JSON validity: Python \texttt{json.loads()} parsing
    \item Tool name: exact string match against expected tool
    \item Arguments: set intersection of required keys
    \item Repeat penalty: -0.30 if \texttt{name} matches any previous turn
\end{itemize}

Inter-rater reliability was not measured because scoring is deterministic. This is a \emph{limitation}: deterministic rubrics can miss semantic correctness that human judges would catch.

\end{document}
